%\begin{center}
%\large \bf \runtitle
%\end{center}
%\vspace{1cm}
\chapter*{\runtitle}

\noindent Compositional controller synthesis avoids state explosion by solving the problem piecewise: at each iteration it composes a subset of the plant, synthesises a partial controller, hides the local events and \textit{minimises} the result before returning it to the system. That minimisation step is what keeps the intermediate models from growing back, and it is also what dominates the cost of the method: it is solved by computing \textit{Weak Synthesis Observational Equivalence} (WSOE) with an $O(nm)$ algorithm.

\noindent In this thesis we propose replacing WSOE with \textit{branching bisimilarity}, for which an $O(m \log n)$ algorithm is available. We show that the replacement is not straightforward: under the natural hypothesis of hiding every local action, branching merges states that WSOE keeps apart, which would amount to over-reducing the model. We prove that restricting the internal actions to the \textit{uncontrollable} local ones yields $\sim_{B} \subseteq \sim_{W}$, and that the inclusion is strict. On that basis we implement the algorithm inside MTSA and integrate it into the compositional flow.

\noindent The evaluation was carried out on 2212 real minimisation instances captured from the synthesis flow itself. Correctness was validated against the reference implementation in mCRL2 and against the output of WSOE. As for performance, the final implementation stays below 1000~ms on the largest instances, where WSOE reaches 7600~ms, and in 83\% of the cases both equivalences produce the same reduction. The cost of the change shows up in memory usage.

\bigskip

\noindent\textbf{Keywords:} Discrete Event Systems, Controller Synthesis, Compositional Synthesis, Minimisation, Branching Bisimilarity, WSOE, Partition Refinement, LTS.
